
%%%%%%
%
% $Autor: Wings $
% $Datum: 2020-01-18 11:15:45Z $
% $Pfad: WuSt/Skript/Produktspezifikation/powerpoint/ImageProcessing.tex $
% $Version: 4620 $
%
%%%%%%

\chapter{Deployment}
\section{Software}
\subsection{Using of TensorFlow}
TensorFlow: General, uses, user manual  
TensorFlow is an Open source deep learning library developed by the Google brain Team in 2015, it is used to create and train machine learning models and can be used by beginners and experts. 
It also provides APIs in Python, C++ and also for some other programming languages.
tensorflow can be used for classic ML uses however the strengths of TensorFlow lie in the development of complex deep learning applications, such as text or image recognition.
in 2019 TensorFLow 2.0 was launched and focused mainly on simplicity and ease of use. 
- Features
AutoDifferentiation, Eager execution, Distribute, Losses, Metrics, TF.nn, Optimizers
- Usage and extensions
TensorFlow is both the core platform and a library that one can use for machine learning purposes, Keras an open-source software library written in Python, acts as an interface for the TensorFlow library to make creating machine learning models easier.

%https://www.tensorflow.org/guide/basics
Applications
%https://datasolut.com/einfuehrung-in-tensorflow/
advantages and disadvantages 
Tensorflow can not be used in mobile phones or embedded devices because of the limited memory, 
to avoid this inconvenience, Google also developed TensorFlow Lite a tool that makes it possible to use already trained models on small memory devices 
\subsection{Using Of TensorFlow Lite}
\subsection{Creating the Application with the Arduino IDE}
\section{Application "Recognizing Gait Patterns}
\subsection{Installation}
\subsection{Configuration}
\subsection{Tests}
\subsection{In the Field}
\section{Monitoring}
\section{Questions/To dos/ Conclusion}
\section{Appendix}
\subsection{Material List}
\subsection{Data sheets}
%\bigskip
%\medskip
%%\input{tikz/Line.tex}

%\input{tikz/node2.tex}
